\section{Conclusion}
We presented \methodname, a post-training KV cache compression method that learns orthonormal low-rank projection bases by directly minimizing decoder-layer output error, rather than standard proxy objectives based on intermediate attention maps.
Across matched KV cache budgets on Llama 3 8B, \methodname consistently improves the performance-memory tradeoff compared to EigenAttention, reducing C4 perplexity by $11.9$ and increasing MMLU accuracy by $5.4$ points at an iso-memory footprint.
Our analysis reveals that better preservation of decoder layer outputs, specifically their directions, is more predictive of end-to-end performance than merely reconstructing intermediate attention quantities.

\section*{Limitations and Future Work}
First, we calibrate key and value bases independently for each layer; explicitly modeling the coupled effect of joint $K$ and $V$ compression and cross-layer error interactions remains an open challenge.
Second, while we validate our approach on Llama 3 with standard sequence lengths, extending the evaluation to longer context settings ($>2048$) and a broader range of architectures is an important next step.
Third, our predictor relies on mean and diagonal variance; employing richer yet efficient summaries (e.g., covariance sketches or low rank moment features) could improve adaptivity to cross-dimension structure.
%Finally, \methodname currently trains predictors for multiple candidate ranks; reducing this overhead through amortized multi-rank training or distillation is a practical direction for deployment.